\documentclass[12pt,a4paper]{article}

\usepackage[utf8]{inputenc}
\usepackage[T1]{fontenc}
\usepackage[margin=2.5cm]{geometry}
\usepackage{hyperref}
\usepackage{microtype}
\usepackage{parskip}
\usepackage{lmodern}

\hypersetup{colorlinks=true, urlcolor=blue!60!black}

\pagestyle{empty}

\begin{document}

\begin{flushright}
Kundai Farai Sachikonye \\
Auf der Lichtung 19, 82178 Puchheim \\
\href{https://github.com/fullscreen-triangle}{github.com/fullscreen-triangle} \\
\texttt{kundai.sachikonye@bitspark.com} \\
(+49) 1626386344 \\[6pt]
15 Mai 2026
\end{flushright}

\vspace{1em}

Bertrandt AG \\
z.\,Hd.\ Tetiana Maksymova \\
Erlangen, 91051

\vspace{1.5em}

\textbf{Bewerbung: AI \& Medical Imaging Engineer (m/w/d) --- Erlangen}

\vspace{1em}

Sehr geehrte Frau Maksymova,

ich bewerbe mich für die Stelle als AI \& Medical Imaging Engineer.
Mein Schwerpunkt liegt auf der Entwicklung von KI-Systemen für biologische und
medizinische Bilddaten sowie auf der Integration großer Sprachmodelle in klinische
Analyse-Workflows --- beides Kernaufgaben dieser Position.

Zu meinen relevantesten Projekten gehören:

\textbf{Medizinische Bildanalyse.}
\href{https://github.com/fullscreen-triangle/helicopter}{Helicopter} ist ein
multimodales Mikroskopie-Analyse-Framework, das Fluoreszenz-, Hellfeld-,
Phasenkontrast- und Spektraldaten durch sequentielle Constraint-Satisfaction
kombiniert und strukturelle Ambiguität aus $\sim$10$^{60}$ möglichen Konfigurationen
auf eindeutige Lösungen reduziert --- ohne Abhängigkeit von einzelnen optischen
Auflösungsgrenzen.
\href{https://github.com/fullscreen-triangle/lavoisier}{Lavoisier} implementiert
eine Dual-Pipeline-Architektur für Massenspektrometrie-Metabolomik: eine numerische
Pipeline und eine visuelle Pipeline, die Spektren in zeitliche Bildsequenzen umwandelt
und durch ein CNN (PyTorch) analysiert. Ergebnis: 98,9\,\% Feature-Extraction-Genauigkeit
gegenüber der NIST-Referenzbibliothek.

\textbf{Generative KI und Agentic AI.}
\href{https://github.com/fullscreen-triangle/zangalewa}{Zangalewa} ist ein
autonomes AI-Code-Workflow-System: Runtime-Exception-Capture → LLM-Diagnose →
Patch-Generierung → git-isolierte Anwendung → Commit oder Rollback. Das System
abstrahiert mehrere LLM-Provider (Claude, OpenAI), verfolgt Token-Nutzung,
implementiert Caching, und enthält eine metakognitive Schicht, die wiederkehrende
Workflows erkennt und Automatisierungsvorschläge generiert. Dies entspricht direkt
dem im Stellenprofil beschriebenen Aufbau von ``GenAI-powered analysis and decision
agents for clinical workflows.'' Die VSCode-Extension
\href{https://github.com/fullscreen-triangle/pugachev-cobra}{Pugachev-Cobra}
bildet das Frontend dieses Systems.

\textbf{Technischer Stack.}
Python ist meine primäre Sprache (PyTorch, NumPy, SciPy, scikit-learn, OpenCV);
dazu TypeScript/Node.js für Backend-Systeme und Rust für Performance-kritische
Komponenten. C\#/.NET habe ich bisher nicht produktiv eingesetzt --- ich möchte
dies transparent kommunizieren. Meine Erfahrung mit stark typisierten Systemen
(TypeScript, Rust) ermöglicht jedoch einen strukturierten Einstieg. DICOM-Integration
und PACS-Anbindung sind neue Domänen für mich; die Datenpipeline-Architektur,
die dafür benötigt wird, entspricht der Arbeit, die ich bereits mit mzML-Datasets
$>$100\,GB und verteiltem Computing (Ray, Dask) geleistet habe.

\textbf{Sprachliches.}
Mein Deutsch ist auf konversationellem Niveau --- ich lebe und arbeite seit 2011
in Deutschland. Technische Dokumentation und schriftliche Kommunikation in Deutsch
kann ich übernehmen; fließendes muttersprachliches Deutsch brächte ich nicht mit.
Englisch ist nativ.

Ich bin in Deutschland arbeitsberechtigt (Aufenthaltserlaubnis) und sofort
verfügbar. Für Rückfragen stehe ich gerne zur Verfügung.

\vspace{1.5em}
Mit freundlichen Grüßen,

\vspace{2.5em}
\textbf{Kundai Farai Sachikonye}

\end{document}
